\section{Why a fast mapper is not yet a database system}
\label{sec:solution-gap-scope}

The previous chapter introduced several ways to search biological
sequences.  A specialized mapper such as minimap2 is already a highly
engineered answer to the algorithmic question.  It avoids whole-matrix
dynamic programming, reduces the seed set with minimizers, chains
compatible anchors, and invokes expensive alignment only for selected
regions.  It also provides a persistent \texttt{.mmi} representation for
reuse of a reference index.  It would therefore be misleading to motivate
BioSeqDB by claiming that existing mappers are generally slow or
algorithmically naive.

The missing capability lies at a different boundary.  A mapper normally
owns a process, reads files or an index image, maps a stream of queries,
and writes another file or stream.  PostgreSQL owns relations, snapshots,
transactions, typed expressions, joins, memory contexts, and a
cost-based plan.  Calling the first system from the second does not make
their execution models coincide.  The database still needs to know when
alignment will run, which rows supply its operands, whether one input row
can yield several outputs, which intermediate data can be reused, and how
the work is bounded when a relation contains billions of bases.

BioSeqDB addresses this mismatch.  Its problem is not merely
``alignment from SQL''.  A scalar PostgreSQL function can invoke a mapper
and satisfy that literal description.  The stronger objective is to make
sequence search a planned, batch-oriented, resource-accounted database
operation while retaining a credible minimap2-derived algorithmic
baseline.  This chapter separates the scientific limitations of
alignment heuristics from the systems limitations created by storage and
query integration.  This separation is important because the remedies
are different.  A missed seed needs an algorithmic remedy.  Repeated
detosting, cold sidecar reads, or one function call per tuple needs a
systems remedy.

\section{Scientific limitations do not disappear with optimization}
\label{sec:scientific-solution-gaps}

\subsection{Seed selection trades work for evidence}

Minimizers make large-scale mapping feasible by retaining a sparse subset
of the available \(k\)-mers.  The reduction is valuable precisely because
the aligner no longer examines every possible local match.  It also means
that the seed set is evidence for an alignment, not a proof that every
biologically meaningful alignment has been retained.  Sensitivity depends
on \(k\), the minimizer window, the hash and tie policy, the treatment of
ambiguous symbols, and the error pattern in the query.  A short query may
contain too few representatives.  A repetitive region may produce a very
large posting list.  A mutation or sequencing error can remove a seed
that would otherwise connect two parts of a chain.

No database integration can eliminate this trade-off.  BioSeqDB must
preserve it and make the selected policy explicit.  A persisted index is
usable only when its alphabet, \(k\), window, layout, hash semantics, and
source generation agree with the query path.  Reading an incompatible
index quickly would be a correctness failure.  Recomputing every index
for safety would discard the central benefit of persistence.  The system
therefore needs compatibility metadata and a fail-closed dispatch policy,
not only serialized posting lists.

The same point applies to vectorized minimizer extraction.  SIMD changes
how several candidate positions are evaluated.  It must not silently
change which positions are emitted unless the alternative sketch is
deliberately defined as a different algorithm.  An early BioSeqDB
incident illustrates the risk.  A circular-buffer slot equality was
interpreted as distance zero even though it represented one complete
window.  The resulting position shift then caused a later valid
minimizer to be removed by position deduplication.  Adjacent off-by-one
and inactive-lane problems were exposed by the same test
campaign.\footnote{\texttt{DOC0068}, ``Why kmers\_test.cpp Was
Failing'', 2026-04-23.}  The lesson is broader than that defect.  An
implementation is not semantically equivalent merely because it is a
vector form of the same pseudocode.

\subsection{Chaining and extension remain heuristic filters}

Chaining ranks collinear anchors without evaluating every possible
alignment path.  Parameters such as maximum predecessor distance,
occurrence filtering, chain score, skip limits, and gap penalties
determine which evidence survives.  Their useful values vary with read
length, error profile, repeat structure, and reference scale.  A policy
that is effective for a long noisy read can remove all evidence for a
short exact probe.  A permissive policy can retain too many candidates
and transfer the cost to sorting, chaining, or final alignment.

The final extension step is exact only within the regions and score model
that reach it.  Affine-gap dynamic programming can establish the best
alignment inside a selected band, but it cannot recover a candidate that
seeding or chaining removed.  Banded or block-based methods introduce
additional resource and boundary choices.  Traceback and CIGAR
construction can also require substantially more memory than score-only
evaluation.  Consequently, ``alignment accuracy'' is not a single
counter.  At minimum, evaluation must distinguish whether a result was
found, whether its coordinate is correct within a stated tolerance,
whether the expected strand was selected, and whether a requested CIGAR
was produced correctly.

These properties explain why BioSeqDB keeps minimap2 as a comparison
baseline and tests intermediate algorithms against biological result
receipts.  A faster stage is not a contribution when it moves an anchor,
changes a primary result, or falls back to another implementation without
reporting the fallback.  Later chapters give the exact policies used for
minimizer extraction, chaining, and extension.

\subsection{No processor path dominates every workload}

Scalar, SIMD, and GPU execution have different fixed and variable costs.
SIMD is attractive when independent lanes can perform similar work with
regular memory access.  It is less attractive when lane-local rescans,
divergent control flow, gathers, or conversion between array layouts
dominates the useful arithmetic.  A GPU adds command submission, buffer
management, synchronization, and transfer or residency requirements.  A
large resident batch can amortize these costs.  One short query against
one small target often cannot.

The current project reports demonstrate both sides of this boundary.  A
resident multi-query GPU path can enter GPU seeding, chaining, and bounded
CIGAR construction with no CPU or SIMD fallback.  The same reports state
that the general G3SA traceback design is not reproduced in every case
and that a small bounded-CIGAR workload remained faster in minimap2.
Likewise, a one-query multi-gigabase workload is not valid evidence for a
resident GPU batch when the batch branch is not admitted.\footnote{
\texttt{DOC0356}, batched resident GPU implementation report, and
\texttt{DOC0369}, local end-to-end benchmark report dated 2026-07-08.}

BioSeqDB therefore needs an explicit dispatch contract.  It must record
which path was selected, why an accelerator declined work, whether a
fallback occurred, and which correctness receipts belong to the selected
path.  ``GPU enabled'' or ``SIMD case'' is not sufficient evidence.
This is an algorithmic and experimental limitation of heterogeneous
execution, not a PostgreSQL-specific defect.

\section{The file-oriented workflow leaves systems work outside the mapper}
\label{sec:file-workflow-gap}

In a conventional workflow, relational metadata and biological sequences
often meet only in orchestration code.  A query selects samples or
identifiers.  A script exports sequences to FASTA or FASTQ.  A mapper
opens or constructs an index and emits PAF, SAM, or BAM.  Another step
parses the results and joins them with annotations.  Workflow engines can
make this process repeatable, but the mapper and the database still
optimize different fragments of the computation.

This separation has several consequences.

\begin{enumerate}
  \item The database planner cannot compare an alignment path with another
  relational plan because the expensive operation occurs after export.
  Selective predicates and joins may be applied too early, too late, or in
  duplicated scripts.

  \item Sequence bytes cross representation boundaries.  They may be
  decoded from a packed database value, rendered as text, parsed by the
  mapper, and parsed again when results return.

  \item The lifetime of an index is external to the source relation.
  Scripts must decide whether an index is current, compatible, and safe to
  reuse after updates.

  \item Failures span several durability domains.  A committed database
  change, a partially written file, and a terminated mapper process do not
  automatically form one atomic operation.

  \item Alignment results are detached from SQL expression semantics.
  Producing several matches for one sequence requires an explicit result
  table or file expansion before relational processing can resume.
\end{enumerate}

These costs are not evidence that file formats are unnecessary.  FASTA,
FASTQ, PAF, SAM, BAM, and mapper indexes are valuable interchange and
archival formats.  The gap is that they become mandatory execution
boundaries even when both the source data and the desired analysis already
belong to one database query.

\section{Why an ordinary SQL function is insufficient}
\label{sec:scalar-function-gap}

PostgreSQL makes it straightforward to expose C or C++ code as a typed
function.  A direct implementation of
\(\operatorname{mmsearch}(target, query)\) could decode its two arguments,
run the mapper, and return a Boolean or a result value.  This form is a
useful correctness fallback.  It is a poor primary execution model for
BioSeqDB.

First, PostgreSQL normally evaluates a scalar expression for one tuple at
a time.  That makes batching accidental.  It prevents a scan from
collecting many targets or queries before choosing an implementation, and
it makes accelerator submission costs recur at tuple granularity.
Second, one sequence pair can produce zero, one, or several matches.  A
Boolean predicate alone cannot express the result cardinality needed by
the projection.  Third, score, coordinates, strand, CIGAR, divergence,
and mapping quality belong to the same alignment result.  Re-evaluating
one expensive function for each projected field would duplicate work.
Fourth, the ordinary function is evaluated before a later batch could
know that a row has no result.

The early custom-scan proof exposed these issues directly.  It could
replace one function call with one mutable constant and could buffer
tuples, but it did not yet invoke the real aligner, return the complete
result, filter a no-match row after batch execution, support the
vectorized scan path, or emit multiple rows for one input.\footnote{
\texttt{DOC0054}, ``Wiring the real aligner batch through the Postgres
custom scan'', 2026-04-21.  This archived design is used as history.
The current contract is stated in \texttt{DOC0321}.}

The present design keeps an ordinary function body so that PostgreSQL can
resolve types and so that unsupported contexts have a defined fallback.
Its planner cost is intentionally high.  A support function and relation
path hook offer a cheaper custom scan when the expression has a supported
shape.  The custom plan owns batching and replaces result-bearing
expressions with nodes whose values are populated for each emitted match.
The custom scan can therefore return no row for a target without a match
or several rows for a target with several matches.  Current project
guidelines require this plan choice to be verified with
\texttt{EXPLAIN (ANALYZE)} and require both vectorized and
non-vectorized execution to be tested.\footnote{\texttt{DOC0321},
Sections 7.2 to 7.4, current code and workflow guidelines.}

This planner arrangement also creates obligations that a standalone
mapper does not have.  PostgreSQL may copy expression trees, rescan an
inner plan, evaluate constants during planning, or reset a memory context.
The extension must keep node identity, ownership, and rescan state valid
under those operations.  Plan costs and EXPLAIN annotations must describe
the path that will actually execute.  Hidden thread-local state is
inappropriate because it changes function behavior without changing its
arguments and interacts badly with constant folding and error unwinding.

\section{Representation and lifetime are part of performance}
\label{sec:representation-gap}

A biological sequence is logically a range of symbols.  At the storage
boundary it may be a PostgreSQL varlena value, an inline packed payload, a
fragmented column, a lazy slice, a concatenation, a reverse view, a
mapped external page, or a materialized C++ container.  Treating all of
these as \texttt{std::string} would simplify an interface by hiding the
most important cost.

The April integration design illustrates the problem.  A Datum obtained
from a tuple slot cannot be retained after the slot is reused.  A toasted
value may have to be detoasted into memory owned by a PostgreSQL context.
The first real batch design therefore copied each complete varlena into an
owned C++ byte vector and then constructed an owning
\texttt{compact\_sequence\_range} for each row.  This was a reasonable
correctness step.  The same document explicitly listed a non-owning
compact view and batch-level target-index reuse as follow-up
work.\footnote{\texttt{DOC0058}, Sections 4, 7, and 11, archived
2026-04-21 design.}

At scale, ownership conversion can dominate an algorithm that was
optimized to avoid touching most symbols.  Eagerly materializing a long
fragmented sequence creates work proportional to the whole sequence even
when minimizer extraction or alignment needs one subrange.  Reversing by
copying doubles traffic.  Decoding a packed sequence to bytes discards
the storage density that made database scans feasible.  Rebuilding a
contiguous target so worker threads can share it may be worthwhile, but
that cost must be a deliberate stage with a memory budget rather than an
implicit function argument conversion.

This motivates the common sequence abstraction developed later in the
thesis.  Algorithms consume a compact range contract and convert only the
needed region at the point of use.  PostgreSQL-specific ownership remains
outside the numeric kernels.  Packed, fragmented, and lazy views must
preserve alphabet information, base-pair length, and iteration semantics.
Zero copy is not an unconditional objective.  It is useful only when the
borrowed lifetime is proven and when irregular access does not cost more
than one bounded materialization.

\section{A persisted index is a database object, not merely a byte image}
\label{sec:persistence-gap}

The ability to save an index is necessary but not sufficient for
database-resident reuse.  A PostgreSQL index representation must answer
questions that an external reference image can leave to the caller:

\begin{enumerate}
  \item Which exact source relation and generation produced the index?

  \item Which alphabet, minimizer parameters, hash scheme, position
  width, and layout can consume it?

  \item Does a transaction see the old generation or the new generation
  while an update is being published?

  \item How are an aborted publication, a dropped table, a tablespace
  move, a crash, and an orphaned generation handled?

  \item Can a cold reader select only the needed extents before allocating
  memory, or must it reconstruct the entire index?

  \item How are checksum work, decoded buffers, posting lists, GPU
  residency, and temporary query projections charged to one memory
  budget?
\end{enumerate}

The project's first answer was a shared sidecar keyed by relation and row
number with a flat serialized target index.  It was deliberately a
proof-of-concept.  It accepted orphan cleanup, native-layout serialization,
and coarse shared-table locality.  Later designs moved toward immutable
generations, explicit wire formats, compatibility manifests, bounded
selective loading, and PostgreSQL-owned publication.  This progression is
important.  It shows why ``store the minimizers in a bytea'' describes the
first write operation but not the lifecycle of an index.

Cold loading also changes the performance question.  A precomputed index
can be slower than recomputation when the reader performs many small tuple
fetches, validates the same payload twice, builds a query-projected facade,
sorts data again, or repeatedly walks all selected extents for memory
accounting.  The July 15 diagnosis found both an
\(O(E^2)\) ownership-accounting path over selected extents and a duplicate
4.46 GB checksum traversal in the then-current Format 5 reader.
It explicitly warned that its payload timer also included C++ routing and
decode work, so it could not be interpreted as a pure PostgreSQL storage
timer.\footnote{\texttt{DOC0381}, Sections 2 and 3,
``Bioseq just\_index cold-load acceleration'', 2026-07-15.}

The resulting requirement is stronger than persistence.  The stored form
must already support the dominant query access pattern, and the reader
must authenticate and materialize it once under a linear memory ledger.
The index chapter describes the current generation and loading design.

\section{Bounded memory and parallel work must be designed together}
\label{sec:resource-gap}

Sequence workloads combine large immutable inputs with variable-size
derived structures.  Target minimizers, posting lists, query hashes,
anchors, chain state, alignment matrices, traceback, and output rows can
be live at overlapping times.  A planner estimate based only on source
sequence bytes misses this amplification.  A GPU path adds device buffers
and often a host staging representation.

The naive responses are both undesirable.  Keeping every intermediate
resident can exceed the backend or device budget.  Spilling immediately
can leave most CPU cores idle while storage becomes the bottleneck.
BioSeqDB instead needs phase-aware admission.  Before a task starts, the
system estimates its source, derived, scratch, and output ownership.  It
admits parallel tasks while their combined peak remains within the
budget.  Buffers are released or their authority is transferred when a
phase completes.  External runs or recomputation remain fallbacks, not
the default while memory is available.

PostgreSQL adds a thread-affinity constraint.  Catalog access, tuple
formation, snapshots, and many memory-context operations must remain on
the backend thread.  CPU workers can operate on immutable C++ views after
the backend has acquired and validated them.  Their tasks must join
before the backend releases PostgreSQL-owned resources.  This boundary
prevents a direct translation of a standalone mapper's thread pool.
Task parallelism is valuable only after input ownership is safe and the
PostgreSQL calls have been removed from worker closures.

These concerns explain an otherwise surprising observation from the
no-index investigation.  A benchmark labeled parallel did not prove that
PostgreSQL-backed target-index construction was parallel.  The fast cells
benefited from prebuilt indexes and backend-local residency, while the
no-index path included serial loading and failed worker-shareability
gates for lazy PostgreSQL-backed sequences.\footnote{\texttt{DOC0318},
executive verdict and root-cause sections, 2026-06-07.}  Stage labels
must therefore be supported by execution receipts and resource
accounting.

\section{Ingestion, metadata, and external data are part of the query}
\label{sec:ingestion-gap}

An integrated system must load data reliably before it can align it.
FASTA and FASTQ records are not only symbol strings.  They contain
identifiers, descriptions, quality information, source-specific fields,
and sometimes multiline records.  Compression and remote delivery change
how much input can be buffered.  A database loader must turn this stream
into typed columns while preserving record boundaries and reporting the
record that caused a conversion error.

This task should not be embedded in the mapper.  The non-PostgreSQL IO
module parses and decompresses sequence records.  A concept-constrained
writer defines the boundary.  The PostgreSQL implementation maps parsed
fields to destination columns and invokes the appropriate type input
functions.  The current implementation report records a single-pass
FASTA and FASTQ header parser, typed multi-column mapping, integration
tests, and a non-empty \texttt{mmsearch} over data loaded through that
path.\footnote{\texttt{DOC0319}, implementation report dated
2026-06-09.}

Remote sequence databases create the same separation at another scale.
An external provider can expose data as a relation, but network failure,
pagination, authentication, retry policy, and bounded streaming remain
provider responsibilities.  The alignment pipeline should see the same
sequence abstraction whether the source is a local column, a loaded file,
or an external relation.  Otherwise every new source multiplies the
algorithm implementations and their correctness surface.

\section{Evidence from intermediate results}
\label{sec:solution-gap-evidence}

The strongest illustration of the systems gap is that the ranking of two
implementations changed when the index delivery path changed.  Table
\ref{tab:integration-evidence} contains three dated measurements.  They
are not combined into one speedup claim because the reports use different
artifacts and acceptance states.

\begin{table}[ht]
  \centering
  \small
  \begin{tabular}{p{0.16\linewidth}p{0.25\linewidth}p{0.20\linewidth}p{0.29\linewidth}}
    \toprule
    Evidence & BioSeqDB measured region & minimap2 measured region &
    Interpretation \\
    \midrule
    2026-07-08, 2 Gbp DNA4, no index, MPF &
    23.888 s &
    73.566 s &
    Both found the expected result.  The BioSeqDB recompute path was
    faster in this case. \\
    2026-07-08, 2 Gbp DNA4, forced cold restart, MPF persisted index &
    36.187 s &
    15.465 s &
    Both produced the expected match.  BioSeqDB was 2.34 times slower,
    showing that persisted delivery dominated the saved computation. \\
    2026-07-17 handoff, six-run lower-middle P50 &
    5.368746 s &
    16.487527084 s &
    The report records 16 of 16 parity and a 3.071020 times advantage
    after a fragment-native and compact-direct redesign.  Its final
    isolated rerun was still pending. \\
    \bottomrule
  \end{tabular}
  \caption{Intermediate performance evidence for the integration problem.
    The first two rows come from \texttt{DOC0369}, reports 73 and 76.
    The third comes from the current summary of \texttt{DOC0384}.
    The rows demonstrate that storage, reconstruction, and batching can
    reverse the ranking of the same broad alignment pipeline.  They do not
    establish a machine-independent final speedup.}
  \label{tab:integration-evidence}
\end{table}

The first two rows are especially useful because they reject a simple
story.  BioSeqDB could be faster than minimap2 while rebuilding
minimizers, yet slower after being given a persisted index.  The
alignment kernels alone cannot explain that ordering.  The forced restart
removed backend-local warmth, and the later investigation identified
index reconstruction, query projection, checksum, tuple delivery, and
memory-accounting work.  The July 17 handoff then reported that a
fragment-native and compact-direct route reversed the comparison on its
paired framework samples.

The third row is reported with its open gate.  A lower-middle P50 from six
observations is valuable evidence for a design decision.  It is not a
license to claim that every BioSeqDB query is three times faster than
minimap2.  The report itself required a final contemporary rerun and
separated the retained \texttt{just\_index} result from later XXL memory
work.  This thesis follows that distinction.

\section{Comparison of solution classes}
\label{sec:solution-class-comparison}

Table \ref{tab:solution-class-gaps} summarizes the roles of the principal
solution classes.  The wording is deliberately scoped.  A blank in the
BioSeqDB problem column does not mean that a tool is defective.  It means
that the concern belongs to another layer.

\begin{table}[ht]
  \centering
  \small
  \begin{tabular}{p{0.16\linewidth}p{0.22\linewidth}p{0.27\linewidth}p{0.25\linewidth}}
    \toprule
    Solution class & Primary strength & Gap at the BioSeqDB boundary &
    Requirement derived for this thesis \\
    \midrule
    Standalone minimap2 &
    Mature seed, chain, and extend mapping with reusable reference
    indexes &
    Does not own PostgreSQL snapshots, relational joins, tuple
    cardinality, or extension memory contexts &
    Preserve a minimap2-derived baseline while adapting data and control
    flow to a custom scan. \\
    File and workflow pipeline &
    Portable interchange, reproducible orchestration, and composition of
    specialized programs &
    Parsing, index validation, failure handling, and joins remain across
    process and file boundaries &
    Keep file IO but remove mandatory export and re-import for
    database-resident queries. \\
    Scalar PostgreSQL function &
    Simple typed SQL surface with a well-defined fallback &
    Per-row evaluation hides batching, duplicates multi-field work, and
    cannot naturally emit several matches per input &
    Use planner support and a batch-owning custom scan with per-match
    result substitution. \\
    Generic columnar storage &
    Compression, projection, and scan bandwidth for typed columns &
    Does not define sequence alphabets, minimizer compatibility, or
    alignment-specific derived data &
    Add compact sequence types, sequence metadata, and versioned
    persistent indexes. \\
    Accelerator-first mapper &
    High throughput when enough regular work is resident on the device &
    Submission, transfer, memory admission, and fallback can dominate
    small or irregular SQL batches &
    Dispatch by measured workload shape and expose path and fallback
    receipts. \\
    BioSeqDB target design &
    One relational surface for storage, filtering, joining, and batched
    alignment &
    Adds planner, lifecycle, ownership, compatibility, and resource
    obligations to every optimized kernel &
    Treat these obligations as first-class contributions and validate
    them separately from algorithm speed. \\
    \bottomrule
  \end{tabular}
  \caption{Existing solution classes and the additional requirements
    introduced by database-resident alignment.}
  \label{tab:solution-class-gaps}
\end{table}

\section{Problem statement and design requirements}
\label{sec:bioseq-problem-statement}

The design problem can now be stated precisely:

\begin{quote}
  Given biological sequences and relational metadata managed by
  PostgreSQL, execute a minimap2-derived search as a planned,
  set-oriented operation that can reuse compatible persistent indexes,
  operate over compact or lazy representations, remain within explicit
  host and device memory budgets, and return zero or more typed alignment
  rows without sacrificing a correct ordinary SQL fallback.
\end{quote}

This statement yields eight requirements.

\begin{enumerate}
  \item \textbf{Algorithmic fidelity.}  Scalar, SIMD, and GPU variants
  need explicit semantics and correctness receipts.  A dispatch label is
  not proof of the executed algorithm.

  \item \textbf{Planner ownership.}  PostgreSQL must be able to choose a
  custom path whose cost reflects the expensive fallback.  EXPLAIN must
  identify the chosen alignment, index, vectorization, and accelerator
  path at the level needed to audit a query.

  \item \textbf{Set-valued batching.}  The executor must collect enough
  work to amortize indexing and accelerator costs, then emit zero, one, or
  several results for each input row.

  \item \textbf{Representation independence.}  Numeric kernels must
  consume a common compact range instead of PostgreSQL Datum or an eager
  byte vector.  Alphabet and layout compatibility remain explicit.

  \item \textbf{Transactional persistence.}  Derived indexes need a
  versioned format, source identity, atomic publication, invalidation,
  crash and abort behavior, and bounded cold loading.

  \item \textbf{Resource accounting.}  Source, decoded, index, scratch,
  output, and device allocations must belong to one phase-aware budget.
  Parallel tasks may use only resources whose ownership survives their
  lifetime.

  \item \textbf{Reliable data entry.}  Local and remote sources need
  bounded streaming, record-aware errors, typed metadata mapping, and a
  clean boundary between provider code and the mapper.

  \item \textbf{Reproducible evidence.}  Performance results must name the
  workload, index state, cache or restart state, selected implementation,
  fallback counters, biological result receipts, and measured region.
\end{enumerate}

The requirements can conflict.  A zero-copy view can prevent safe worker
sharing.  A persistent representation optimized for dense reading can
waste work for a selective query.  A GPU batch can improve throughput
while increasing latency.  An exact compatibility check can add cold
startup work.  The design must make these trade-offs visible instead of
hiding them behind a single function.

\section{Research questions and contributions}
\label{sec:solution-gap-questions}

The thesis investigates four questions derived from the gap analysis.

\begin{enumerate}
  \item How can a minimap2-derived pipeline operate over packed,
  fragmented, PostgreSQL-backed, and lazily composed sequences without
  converting every input to one owning byte representation?

  \item How can PostgreSQL plan a familiar typed expression while a custom
  executor performs batch-oriented, potentially set-valued alignment and
  supplies all projected result fields from one execution?

  \item How can minimizer indexes and derived metadata be stored,
  versioned, validated, loaded, cached, and replaced under PostgreSQL
  lifecycle rules without making cold reuse slower than recomputation?

  \item Under which workload and resource conditions should scalar, SIMD,
  or Vulkan stages execute, and what evidence is sufficient to show that
  the selected path was both faster and biologically acceptable?
\end{enumerate}

The contributions claimed in this thesis are correspondingly
systems-oriented and algorithm-aware.

\begin{enumerate}
  \item A common compact sequence interface connects owning, bit-packed,
  fragmented, PostgreSQL-backed, and lazy representations to the same
  pipeline while keeping alphabet dispatch explicit.

  \item A custom-scan execution model converts supported SQL expressions
  into batch alignment, supports one-to-many results, and substitutes
  score, coordinate, strand, and CIGAR fields without re-running the
  mapper for every projection.

  \item PostgreSQL-resident sequence and minimizer storage ties derived
  generations to source identity and supports compatibility checks,
  bounded loading, caching, and controlled recomputation.

  \item CPU vector and Vulkan implementations are integrated behind
  explicit selection and fallback contracts rather than exposed as
  unconditional replacements for the baseline.

  \item Bounded ingestion, typed header mapping, external providers, and
  generated PostgreSQL boundary code extend the same typed data plane
  beyond the alignment kernels.
\end{enumerate}

These are not claims that BioSeqDB replaces every mapper, workflow engine,
or sequence format.  They establish a design for workloads whose data,
metadata, and desired results already benefit from relational storage and
query planning.  They also do not imply that every experimental
optimization became part of the final route.  The later chapters identify
which intermediate algorithms were retained, which were rejected, and
which benchmark gates remained open at the documented cutoff.

\section{Scope of the evaluation}
\label{sec:solution-gap-evaluation-scope}

Three cautions govern the evaluation in the remainder of the thesis.
First, minimap2 is a baseline for comparable mapping work, not a proxy for
all external workflow costs.  A comparison must state whether loading,
index construction, index opening, PostgreSQL restart, result parsing, and
correctness audits are inside the timer.  Second, one local machine cannot
establish portability or universal throughput.  The July 8 report, for
example, explicitly identifies its environment as Darwin on Arm and
rejects conclusions about x86-64 Linux.  Third, historical reports describe
the implementation that produced them.  A newer architecture can
supersede their design conclusions without making the earlier measurement
false.

For these reasons, the thesis uses numbers to answer bounded questions.
A microbenchmark identifies a kernel bottleneck.  A selective end-to-end
case demonstrates SQL integration and a real result.  A cold-index run
tests persistence rather than backend-local cache residency.  A GPU
result needs branch, dispatch, and fallback receipts.  No single chart is
asked to prove all of these properties.

The next chapter follows one query through the resulting pipeline.  It
defines the stage boundaries that make these separate forms of evidence
possible.  The algorithm chapters then compare each retained stage with
its minimap2 or research origin.  The systems chapters return to the
sequence, planner, storage, index, and IO requirements established here.
